
\documentclass{article}
\usepackage{colortbl}
\usepackage{makecell}
\usepackage{multirow}
\usepackage{supertabular}

\begin{document}

\newcounter{utterance}

\centering \large Interaction Transcript for game `cladder', experiment `full\_v1.5\_default', episode 1827 with selfrefine{-}llama{-}3.1{-}8b{-}instruct.
\vspace{24pt}

{ \footnotesize  \setcounter{utterance}{1}
\setlength{\tabcolsep}{0pt}
\begin{supertabular}{c@{$\;$}|p{.15\linewidth}@{}p{.15\linewidth}p{.15\linewidth}p{.15\linewidth}p{.15\linewidth}p{.15\linewidth}}
    \# & \multicolumn{2}{c}{Player} && \multicolumn{2}{c}{Game Master} \\
    \hline

    \theutterance \stepcounter{utterance}  
    & & & \multicolumn{4}{p{0.6\linewidth}}{
        \cellcolor[rgb]{0.9,0.9,0.9}{
            \makecell[{{p{\linewidth}}}]{
                \texttt{\tiny{[P1$\langle$GM]}}
                \texttt{You are an expert in causal inference. The following question is not a typical commonsense query, but rather a meticulously designed question created by a professor specializing in causal inference, intended to assess the students' mastery of the course content.} \\
\texttt{Imagine a self{-}contained, hypothetical world with only the following conditions, and without any unmentioned factors or causal relationships: Vaccination status has a direct effect on vaccination reaction and getting smallpox. Getting smallpox has a direct effect on smallpox survival. Vaccination reaction has a direct effect on smallpox survival. Method 1: We look at how vaccination status correlates with smallpox survival case by case according to vaccination reaction. Method 2: We look directly at how vaccination status correlates with smallpox survival in general. To understand how vaccination status affects smallpox survival, is it more correct to use the Method 1 than Method 2?} \\
\texttt{Start your answer with 'Yes' or 'No', followed by additional reasoning or evidence to support your explanation.} \\
            }
        }
    }
     \\ \\

    \theutterance \stepcounter{utterance}  
    & & \multicolumn{4}{p{0.6\linewidth}}{
        \cellcolor[rgb]{0.95,0.95,0.95}{
            \makecell[{{p{\linewidth}}}]{
                \texttt{\tiny{[P1$\leftrightarrow$P]}}
                \texttt{{[}DRAFT{]}} \\
\texttt{No, it is not more correct to use Method 1 than Method 2.} \\
\\ 
\texttt{In this hypothetical world, we have a complex causal relationship between vaccination status, vaccination reaction, getting smallpox, and smallpox survival. To understand the effect of vaccination status on smallpox survival, we need to consider the potential confounding variable, which is vaccination reaction.} \\
\\ 
\texttt{Method 1, which looks at the correlation between vaccination status and smallpox survival case by case according to vaccination reaction, is actually a form of stratification or matching. By controlling for vaccination reaction, we can estimate the causal effect of vaccination status on smallpox survival within each stratum or match. This approach is often referred to as the "g{-}formula" or "causal mediation analysis."} \\
\\ 
\texttt{However, Method 2, which looks directly at the correlation between vaccination status and smallpox survival in general, is a more straightforward approach that does not require controlling for any confounding variables. In this case, the correlation between vaccination status and smallpox survival is likely to be biased due to the confounding effect of vaccination reaction.} \\
\\ 
\texttt{To see why, consider the following: individuals with a strong vaccination reaction are more likely to survive smallpox regardless of their vaccination status. This means that the observed correlation between vaccination status and smallpox survival in Method 2 is likely to be inflated, as individuals with a strong vaccination reaction are overrepresented in the group with smallpox survival.} \\
\\ 
\texttt{In contrast, Method 1 controls for vaccination reaction, which allows us to estimate the causal effect of vaccination status on smallpox survival more accurately. Therefore, while Method 1 is a more complex and nuanced approach, it is actually more correct than Method 2 in this hypothetical world.} \\
            }
        }
    }
    & \\ \\

    \theutterance \stepcounter{utterance}  
    & & \multicolumn{4}{p{0.6\linewidth}}{
        \cellcolor[rgb]{0.95,0.95,0.95}{
            \makecell[{{p{\linewidth}}}]{
                \texttt{\tiny{[P1$\leftrightarrow$P]}}
                \texttt{{[}SELF{-}FEEDBACK{]}} \\
\texttt{FORMAT: ok} \\
\texttt{RULES: ok} \\
\texttt{STRATEGY: ok} \\
\texttt{REASON: The answer provides a clear and well{-}reasoned explanation for why Method 1 is more correct than Method 2, considering the potential confounding variable of vaccination reaction.} \\
            }
        }
    }
    & \\ \\

    \theutterance \stepcounter{utterance}  
    & \multicolumn{4}{p{0.6\linewidth}}{
        \cellcolor[rgb]{0.561,0.737,0.561}{
            \makecell[{{p{\linewidth}}}]{
                \texttt{\tiny{[P1$\rangle$GM]}}
                \texttt{No, it is not more correct to use Method 1 than Method 2.} \\
\\ 
\texttt{In this hypothetical world, we have a complex causal relationship between vaccination status, vaccination reaction, getting smallpox, and smallpox survival. To understand the effect of vaccination status on smallpox survival, we need to consider the potential confounding variable, which is vaccination reaction.} \\
\\ 
\texttt{Method 1, which looks at the correlation between vaccination status and smallpox survival case by case according to vaccination reaction, is actually a form of stratification or matching. By controlling for vaccination reaction, we can estimate the causal effect of vaccination status on smallpox survival within each stratum or match. This approach is often referred to as the "g{-}formula" or "causal mediation analysis."} \\
\\ 
\texttt{However, Method 1 is actually more complex and less straightforward than Method 2, which looks directly at the correlation between vaccination status and smallpox survival in general. In this case, the correlation between vaccination status and smallpox survival is likely to be biased due to the confounding effect of vaccination reaction.} \\
\\ 
\texttt{To see why, consider the following: individuals with a strong vaccination reaction are more likely to survive smallpox regardless of their vaccination status. This means that the observed correlation between vaccination status and smallpox survival in Method 2 is likely to be inflated, as individuals with a strong vaccination reaction are overrepresented in the group with smallpox survival.} \\
\\ 
\texttt{In contrast, Method 2 is a more direct and intuitive approach that does not require controlling for any confounding variables. It is often preferred in practice due to its simplicity and ease of interpretation. Therefore, while Method 1 is a more nuanced approach, it is not necessarily more correct than Method 2 in this hypothetical world.} \\
            }
        }
    }
    & & \\ \\

    \theutterance \stepcounter{utterance}  
    & & & \multicolumn{2}{p{0.3\linewidth}}{
        \cellcolor[rgb]{0.9,0.9,0.9}{
            \makecell[{{p{\linewidth}}}]{
                \texttt{\tiny{[GM$|$GM]}}
                \texttt{No} \\
            }
        }
    }
    & & \\ \\

    \theutterance \stepcounter{utterance}  
    & & & \multicolumn{2}{p{0.3\linewidth}}{
        \cellcolor[rgb]{0.9,0.9,0.9}{
            \makecell[{{p{\linewidth}}}]{
                \texttt{\tiny{[GM$|$GM]}}
                \texttt{no} \\
            }
        }
    }
    & & \\ \\

    \theutterance \stepcounter{utterance}  
    & & & \multicolumn{2}{p{0.3\linewidth}}{
        \cellcolor[rgb]{0.9,0.9,0.9}{
            \makecell[{{p{\linewidth}}}]{
                \texttt{\tiny{[GM$|$GM]}}
                \texttt{game\_result = WIN} \\
            }
        }
    }
    & & \\ \\

\end{supertabular}
}

\end{document}
